\documentclass[11pt,a4paper]{scrartcl}

\usepackage{amssymb}
\usepackage{amsmath}

\usepackage[utf8]{inputenc}
\usepackage[T1]{fontenc}
\newcommand{\ats}[1]{\par\noindent\textit{Atsakymas:} #1}
\usepackage{cancel}
\usepackage{tikz}
\usepackage{lmodern} 
\usepackage{amsmath,amssymb,amsthm}
\usepackage{graphicx}
\usepackage[dvipsnames,svgnames]{xcolor}
\usepackage{hyperref}
\usepackage{mathrsfs}
\usepackage[shortlabels]{enumitem}
\usepackage[obeyFinal,textsize=scriptsize,shadow,loadshadowlibrary]{todonotes}
\usepackage{mathtools}
\usepackage{microtype}

%%% Paragraph layout
\setlength{\parindent}{0pt}
\setlength{\parskip}{0.6\baselineskip}

%%% Colored sections
\renewcommand*{\sectionformat}%
  {\color{purple}\S\thesection\enskip}
\renewcommand*{\subsectionformat}%
  {\color{purple}\S\thesubsection\enskip}
\renewcommand*{\subsubsectionformat}%
  {\color{purple}\S\thesubsubsection\enskip}
\addtokomafont{paragraph}{\color{orange!35!black}\P\ }
\KOMAoptions{numbers=noenddot}

%%% Title
\addtokomafont{subtitle}{\Large}
\setkomafont{author}{\Large\scshape}
\setkomafont{date}{\Large\normalsize}
% Neigiamas vertikalus tarpas, kuris pakelia dokumento antraštę į viršų. 
% Galite keisti -2.5cm į -3cm (aukščiau) arba -1.5cm (žemiau).
\titlehead{\vspace*{-7.5cm}} 

% PRIDĖTA: Automatiškai perrašome datą į mokyklos pavadinimą
\AtBeginDocument{%
  \date{\normalsize Vilniaus licėjus}
}

%%% Page setup
\usepackage[headsepline]{scrlayer-scrpage}
\renewcommand{\headfont}{}
\addtolength{\textheight}{3.14cm}
\setlength{\footskip}{0.5in}
\setlength{\headsep}{10pt}
% Puslapio antraštė turi rodyti tik konkrečius dokumento duomenis.
% Nustatome juos aiškiai, kad neatsirastų "\@author" / "\@title" arba senų reikšmių.
\makeatletter
\AtBeginDocument{%
  \ihead{\footnotesize\textbf{Andrius Berniukevičius}}%
  \ohead{\footnotesize\textbf{\@title}}%
}
\makeatother
\automark{section}
\chead{}
\cfoot{\pagemark}

%%% Macros
\providecommand{\ol}{\overline}
\providecommand{\ul}{\underline}
\providecommand{\wt}{\widetilde}
\providecommand{\wh}{\widehat}
\providecommand{\eps}{\varepsilon}
\providecommand{\half}{\frac{1}{2}}
\providecommand{\inv}{^{-1}}
\newcommand{\dang}{\measuredangle} %% Directed angle
\newcommand{\vocab}[1]{\textbf{\color{blue}\sffamily #1}}
\providecommand{\CC}{\mathbb C}
\providecommand{\FF}{\mathbb F}
\providecommand{\NN}{\mathbb N}
\providecommand{\QQ}{\mathbb Q}
\providecommand{\RR}{\mathbb R}
\providecommand{\ZZ}{\mathbb Z}
\providecommand{\ts}{\textsuperscript}
\providecommand{\dg}{^\circ}
\providecommand{\ii}{\item}
\providecommand{\defeq}{\coloneqq}
\DeclareMathOperator*{\lcm}{lcm}
\DeclareMathOperator*{\argmin}{arg min}
\DeclareMathOperator*{\argmax}{arg max}
\providecommand{\hrulebar}{\par
\hspace{\fill}\rule{0.95\linewidth}{.7pt}\hspace{\fill}
\par\nointerlineskip \vspace{\baselineskip}}

%%% Optional colored theorems
\usepackage{thmtools}
\usepackage[framemethod=TikZ]{mdframed}
\mdfdefinestyle{mdbluebox}{%
  roundcorner=10pt,
  linewidth=1pt,
  skipabove=12pt,
  innerbottommargin=9pt,
  skipbelow=2pt,
  linecolor=blue,
  nobreak=true,
  backgroundcolor=TealBlue!5,
}
\declaretheoremstyle[
  headfont=\sffamily\bfseries\color{MidnightBlue},
  mdframed={style=mdbluebox},
  headpunct={\\[3pt]},
  postheadspace={0pt}
]{thmbluebox}
\mdfdefinestyle{mdredbox}{%
  linewidth=0.5pt,
  skipabove=12pt,
  frametitleaboveskip=5pt,
  frametitlebelowskip=0pt,
  skipbelow=2pt,
  frametitlefont=\bfseries,
  innertopmargin=4pt,
  innerbottommargin=8pt,
  nobreak=true,
  backgroundcolor=Salmon!5,
  linecolor=RawSienna,
}
\declaretheoremstyle[
  headfont=\bfseries\color{RawSienna},
  mdframed={style=mdredbox},
  headpunct={\\[3pt]},
  postheadspace={0pt},
]{thmredbox}
\mdfdefinestyle{mdgreenbox}{%
  skipabove=8pt,
  linewidth=2pt,
  rightline=false,
  leftline=true,
  topline=false,
  bottomline=false,
  linecolor=ForestGreen,
  backgroundcolor=ForestGreen!5,
}
\declaretheoremstyle[
  headfont=\bfseries\sffamily\color{ForestGreen!70!black},
  bodyfont=\normalfont,
  spaceabove=2pt,
  spacebelow=1pt,
  mdframed={style=mdgreenbox},
  headpunct={ --- },
]{thmgreenbox}
\mdfdefinestyle{mdblackbox}{%
  skipabove=8pt,
  linewidth=3pt,
  rightline=false,
  leftline=true,
  topline=false,
  bottomline=false,
  linecolor=black,
  backgroundcolor=RedViolet!5!gray!5,
}
\declaretheoremstyle[
  headfont=\bfseries,
  bodyfont=\normalfont\small,
  spaceabove=0pt,
  spacebelow=0pt,
  mdframed={style=mdblackbox}
]{thmblackbox}

%% Aplinkos su rėmeliais (thmtools)
\declaretheorem[style=thmbluebox,name=Teorema]{theorem}
\declaretheorem[style=thmbluebox,name=Lema]{lemma}
\declaretheorem[style=thmbluebox,name=Savybė]{property}
\declaretheorem[style=thmbluebox,name=Teiginys]{proposition}
\declaretheorem[style=thmbluebox,name=Išvada]{corollary}
\declaretheorem[style=thmbluebox,name=Teorema,numbered=no]{theorem*}
\declaretheorem[style=thmbluebox,name=Lema,numbered=no]{lemma*}
\declaretheorem[style=thmbluebox,name=Teiginys,numbered=no]{proposition*}
\declaretheorem[style=thmbluebox,name=Išvada,numbered=no]{corollary*}
\declaretheorem[style=thmgreenbox,name=Algoritmas]{algorithm}
\declaretheorem[style=thmgreenbox,name=Algoritmas,numbered=no]{algorithm*}
\declaretheorem[style=thmgreenbox,name=Teiginys]{claim}
\declaretheorem[style=thmgreenbox,name=Teiginys,numbered=no]{claim*}
\declaretheorem[style=thmredbox,name=Pavyzdys]{example}
\declaretheorem[style=thmredbox,name=Pavyzdys,numbered=no]{example*}
\declaretheorem[style=thmblackbox,name=Pastaba]{remark}
\declaretheorem[style=thmblackbox,name=Pastaba,numbered=no]{remark*}

%% Standartinės amsthm aplinkos (Apibrėžimai ir užduotys)
\theoremstyle{definition}
\ifcsname conjecture\endcsname\else\newtheorem{conjecture}{Hipotezė}\fi
\ifcsname definition\endcsname\else\newtheorem{definition}{Apibrėžimas}\fi
\ifcsname fact\endcsname\else\newtheorem{fact}{Faktas}\fi
\ifcsname answer\endcsname\else\newtheorem{answer}{Atsakymas}\fi
\ifcsname ques\endcsname\else\newtheorem{ques}{Klausimas}\fi
\ifcsname exercise\endcsname\else\newtheorem{exercise}{Užduotis}\fi
\ifcsname problem\endcsname\else\newtheorem{problem}{Uždavinys}[section]\fi
\ifcsname conjecture*\endcsname\else\newtheorem*{conjecture*}{Hipotezė}\fi
\ifcsname definition*\endcsname\else\newtheorem*{definition*}{Apibrėžimas}\fi
\ifcsname fact*\endcsname\else\newtheorem*{fact*}{Faktas}\fi
\ifcsname answer*\endcsname\else\newtheorem*{answer*}{Atsakymas}\fi
\ifcsname ques*\endcsname\else\newtheorem*{ques*}{Klausimas}\fi
\ifcsname exercise*\endcsname\else\newtheorem*{exercise*}{Užduotis}\fi
\ifcsname problem*\endcsname\else\newtheorem*{problem*}{Uždavinys}\fi

\newcounter{answerssection}
\newcommand{\answerssection}{%
  \setcounter{answerssection}{\value{section}}%
  \section{Atsakymai}%
  \setcounter{problem}{0}%
  \renewcommand{\theproblem}{\arabic{answerssection}.\arabic{problem}}%
}

%% GALUTINIS SPRENDIMAS PROOF APLINKAI
%% --- PRADŽIA: Įrodymo aplinkos sutvarkymas ---
\makeatletter

% 1. Užtikriname, kad žodis būtų "Įrodymas", net jei "babel" paketas bando jį perrašyti
\AtBeginDocument{%
  \renewcommand{\proofname}{Įrodymas}%
  \@ifpackageloaded{babel}{%
    \addto\captionslithuanian{\renewcommand{\proofname}{Įrodymas}}%
  }{}%
}

% 2. Priverstinai pašaliname scrartcl klasės bold šriftą ir paliekame TIK italic
% 3. Sujungiame su tuščiu kvadratėliu dešiniajame kampe.
\renewcommand{\qedsymbol}{\ensuremath{\Box}}
\renewenvironment{proof}[1][\proofname]{\par
  \pushQED{\qed}%
  \normalfont \topsep6\p@\@plus6\p@\relax
  \trivlist
  \item[\hskip\labelsep
        \normalfont\itshape % Priverčia naudoti tik pasvirą šriftą, kaip nuotraukoje
    #1\@addpunct{.}]\ignorespaces
}{%
  \popQED\endtrivlist\@endpefalse
}
\newenvironment{solution}{\par
  \normalfont \topsep6\p@\@plus6\p@\relax
  \trivlist
  \item[\hskip\labelsep
        \normalfont\itshape
    \textit{Sprendimas}\@addpunct{.}]\ignorespaces
}{%
  \endtrivlist\@endpefalse
}
\newcommand{\ans}[1]{\par\noindent\textit{Atsakymas:} #1}
\makeatother


\title{Skaičių teorija: Pirminiai skaičiai ir pilnieji kvadratai}
\author{Andrius Berniukevičius}

\newenvironment{solution}
    {\begin{list}{}{\setlength{\leftmargin}{10pt}\setlength{\rightmargin}{10pt}}\item[]}
    {\end{list}}

\begin{document}

\maketitle

\section{Kai algebra susitinka su skaičių teorija}

Iki šiol mes skaičius dalinome, skaidėme pirminiais dauginamaisiais ir ieškojome jų bendrų daliklių. Tačiau tikroji olimpiadinė magija prasideda tada, kai į pagalbą skaičių teorijai ateina \textbf{algebra}. 

Pagalvokite: kuo ypatingas pirminis skaičius $p$? Skirtingai nei sudėtiniai skaičiai, kurie gali būti išskaidyti į daugybę skirtingų kombinacijų ($24 = 6 \cdot 4 = 8 \cdot 3 = 12 \cdot 2$), pirminis skaičius yra visiškai „kietas“. Jį dviejų natūraliųjų skaičių sandauga galime užrašyti \textbf{tik vieninteliu būdu}:
$$ p = 1 \cdot p $$

Šis, atrodytų, banalus faktas tampa neįtikėtinai galingu ginklu, kai mes jį sujungiame su pilnaisiais kvadratais.

\subsection{Baziniai ginklai: Kvadratų skirtumas}

Daugelis iš jūsų mokykloje išmoko greitosios daugybos formulę – kvadratų skirtumą:
$$ a^2 - b^2 = (a - b)(a + b) $$
Mokykloje tai naudojama tiesiog prastinant reiškinius. Olimpiadose tai yra mirtinas ginklas lygtims spręsti. 

Įsivaizduokite, kad turite lygtį: $p = a^2 - b^2$ (kur $p$ – pirminis skaičius, o $a, b$ – natūralieji).
Išskaidykime ją:
$$ p = (a - b) \cdot (a + b) $$
Kairėje pusėje turime pirminį skaičių, o dešinėje – \textit{dviejų} skaičių sandaugą. Bet mes ką tik kalbėjome, kad pirminis skaičius skaidomas tik kaip $1 \cdot p$! Vadinasi, vienas iš šių skliaustų privalo būti lygus $1$, o kitas – $p$. Kadangi $a$ ir $b$ yra natūralieji skaičiai, akivaizdu, kad $(a-b)$ yra mažesnis už $(a+b)$. Todėl mes \textbf{be jokio spėliojimo} gauname griežtą lygčių sistemą:
$$ a - b = 1 $$
$$ a + b = p $$
Šio triuko esmė: \textbf{algebrinis skaidymas verčia mus lyginti dauginamuosius su vienetu ir pačiu pirminiu skaičiumi}.

\subsection{Pilnųjų kvadratų spąstai}
Dar viena svarbi strategija dirbant su kvadratais – „įkalinimas“. Kiekvienas natūralusis skaičius, esantis griežtai tarp dviejų gretimų pilnųjų kvadratų $k^2$ ir $(k+1)^2$, pats \textbf{niekada} negali būti pilnasis kvadratas. (Juk tarp $25$ ir $36$ nėra kito sveikojo skaičiaus, kurį pakėlus kvadratu gautume sveiką skaičių). Jei įrodysite, kad $x$ yra tarp $k^2$ ir $(k+1)^2$, jūs įrodysite, kad $x$ nėra pilnasis kvadratas.

\section{Uždavinių sprendimo pavyzdžiai}

\begin{example}[Tiesioginis kvadratų skirtumas]
Raskite pirminį skaičių $p$, jei žinoma, kad $p = n^2 - 25$ (kur $n$ – natūralusis skaičius).
\end{example}
\begin{solution}
\textbf{Idėjos analizė:} Matome pirminį skaičių, lygų išraiškai, kurią galima išskaidyti. Mūsų tikslas – paversti skirtumą sandauga ir mažiausią dauginamąjį prilyginti 1.

\textbf{Sprendimas:}\\
Taikome kvadratų skirtumo formulę:
$$ p = (n - 5)(n + 5) $$
Kadangi $p$ yra pirminis skaičius, vienas iš dauginamųjų dešinėje privalo būti $1$. Kadangi $n$ yra natūralusis skaičius, $(n-5)$ yra mažesnis už $(n+5)$. Vadinasi:
$$ n - 5 = 1 \implies n = 6 $$
Dabar, kai radome $n$, įstatome jį atgal, kad rastume $p$:
$$ p = (6 - 5)(6 + 5) = 1 \cdot 11 = 11 $$
Būtina patikrinti: ar $11$ yra pirminis skaičius? Taip. 
Atsakymas: $p = 11$. \hfill $\square$
\end{solution}

\begin{example}[Kintamųjų perkėlimas]
Raskite visus pirminius skaičius $p$, kuriems esant reiškinys $7p + 1$ yra pilnasis kvadratas.
\end{example}
\begin{solution}
\textbf{Idėjos analizė:} Sąlyga sako, kad $7p + 1 = k^2$ (kažkokiam natūraliajam $k$). Skirtingai nei pirmame pavyzdyje, čia kvadratų skirtumo tiesiogiai nematome. Bet mes galime jį \textbf{sukurti}, perkėlę vienetą į kitą pusę!

\textbf{Sprendimas:}\\
Užrašome sąlygą lygtimi:
$$ 7p + 1 = k^2 $$
Perkeliame $1$ į dešinę pusę:
$$ 7p = k^2 - 1 $$
Pritaikome kvadratų skirtumą:
$$ 7 \cdot p = (k - 1)(k + 1) $$
Čia mes turime situaciją, kai dviejų \textit{pirminių} skaičių (7 ir $p$) sandauga lygi dviejų iš eilės einančių lyginių/nelyginių skaičių sandaugai. Skaičius 7 privalo dalinti arba $(k-1)$, arba $(k+1)$.
Išnagrinėkime mažus $k$ variantus, nes skirtumas tarp skliaustų tėra 2.
Galimi atvejai, kai vienas skliaustas lygus 7:
1) Tarkime, $k - 1 = 7 \implies k = 8$. Tuomet antras skliaustas $k+1 = 9$. 
Gauname $7 \cdot p = 7 \cdot 9 \implies p = 9$. Bet $9$ nėra pirminis skaičius! Šis variantas netinka.
2) Tarkime, $k + 1 = 7 \implies k = 6$. Tuomet pirmas skliaustas $k-1 = 5$.
Gauname $7 \cdot p = 5 \cdot 7 \implies p = 5$. Patikriname: 5 yra pirminis skaičius. $7 \cdot 5 + 1 = 36 = 6^2$. Tinka!
Atsakymas: $p = 5$. \hfill $\square$
\end{solution}

\begin{example}[Kvadratų spąstai ir išskyrimas]
Įrodykite, kad bet kokiam natūraliajam skaičiui $n > 1$, skaičius $n^4 + 4$ niekada nebus pirminis.
\end{example}
\begin{solution}
\textbf{Idėjos analizė:} Reiškinys $n^4 + 4$ yra suma, o ne skirtumas, todėl įprasta formulė neveikia. Tai labai garsi Sophie Germain tapatybė. Jos pagrindinė idėja – \textbf{pilnojo kvadrato išskyrimas}. Mes dirbtinai pridėsime narį, kad gautume kvadratą, ir tuomet atimsime jį atgal, sukurdami mums taip reikalingą kvadratų skirtumą.

\textbf{Sprendimas:}\\
Prisiminkime formulę $(a+b)^2 = a^2 + 2ab + b^2$. Mūsų turimas $n^4 = (n^2)^2$, o $4 = 2^2$. Trūksta vidurinio nario $2 \cdot n^2 \cdot 2 = 4n^2$. Pridėkime jį ir iškart atimkime:
$$ n^4 + 4 = n^4 + \mathbf{4n^2} + 4 - \mathbf{4n^2} $$
Sujungiame pirmuosius tris narius į pilnąjį kvadratą:
$$ = (n^2 + 2)^2 - 4n^2 $$
Dabar $4n^2$ mes galime užrašyti kaip $(2n)^2$. Gavome nuostabų kvadratų skirtumą!
$$ = (n^2 + 2)^2 - (2n)^2 = (n^2 + 2 - 2n)(n^2 + 2 + 2n) $$
Sutvarkome abėcėlės tvarka:
$$ = (n^2 - 2n + 2)(n^2 + 2n + 2) $$
Kad pradinis skaičius būtų pirminis, vienas iš šių skliaustų privalo būti lygus $1$. 
Mažesnysis skliaustas yra $n^2 - 2n + 2$. Išskirkime jame vėl pilnąjį kvadratą:
$$ n^2 - 2n + 2 = (n - 1)^2 + 1 $$
Kadangi uždavinio sąlyga teigia, kad $n > 1$, tai $(n-1)^2 \ge 1$. Todėl mažiausia mažesniojo skliausto reikšmė bus $1^2 + 1 = 2$.
Nei vienas skliaustas negali būti lygus $1$, o tai reiškia, kad $n^4+4$ visada išsiskleis į dviejų didesnių už vienetą skaičių sandaugą. Todėl jis \textbf{niekada nėra pirminis}. \hfill $\square$
\end{solution}

\vspace{0.5cm}

\newpage
\section{Uždaviniai savarankiškam sprendimui (30 iššūkių)}

Pirmieji 10 uždavinių skirti bazinės technikos įtvirtinimui. Visi sprendimai turi būti argumentuoti lygtimis (skaidymu), o ne atsitiktiniu spėliojimu.

\subsection*{Baziniai (Rankos atmušimas)}
\begin{enumerate}
    \item Raskite pirminį skaičių $p$, jei žinoma, kad $p = n^2 - 1$ (kur $n$ – natūralusis skaičius).
    \item Raskite visus pirminius skaičius $p$, kuriems esant teisinga lygybė $p = n^2 - 9$ (kur $n$ – natūralusis skaičius).
    \item Išspręskite lygtį natūraliaisiais skaičiais: $x^2 - y^2 = 13$. (Prisiminkite, kad 13 – pirminis).
    \item Raskite visus pirminius $p$ ir $q$, kuriems teisinga lygybė $p^2 - q^2 = 17$.
    \item Natūralusis skaičius $N$ turi lygiai $3$ daliklius. Įrodykite, kad šis skaičius privalo būti pirminio skaičiaus kvadratas ($N = p^2$).
    \item Skaičiaus kanoninis skaidinys yra $A = 2^x \cdot 3^4 \cdot 5^y$. Kokios turi būti kintamųjų $x$ ir $y$ savybės (lyginumas), kad $A$ būtų pilnasis kvadratas?
    \item Ar egzistuoja toks pirminis skaičius $p$, kuriam esant $p^2 - 4$ yra taip pat pirminis skaičius?
    \item Įrodykite, kad tarp dviejų gretimų pilnųjų kvadratų $k^2$ ir $(k+1)^2$ nėra jokio kito pilnojo kvadrato. Remiantis tuo, ar gali skaičius $k^2 + 1$ (kai $k \ge 1$) būti pilnasis kvadratas?
    \item Raskite pirminį skaičių $p$, jei žinoma, kad $p = x^2 - y^2$, o $x$ ir $y$ skiriasi lygiai per $1$.
    \item Ar gali pirminis skaičius $p > 2$ būti išreikštas dviejų lyginių skaičių kvadratų skirtumu $x^2 - y^2$? Atsakymą pagrįskite.
\end{enumerate}

\subsection*{Vidutinio sunkumo (Rajoninių olimpiadų lygis)}
\setcounter{enumi}{10}
\begin{enumerate}
    \item Raskite visus pirminius $p$, kuriems esant reiškinys $p^2 + 11$ yra pilnasis kvadratas. (Tarkite, kad $p^2 + 11 = k^2$).
    \item Raskite visus pirminius $p$, su kuriais skaičius $13p + 1$ yra pilnasis kvadratas.
    \item Įrodykite, kad joks natūralusis skaičius pavidalo $n^2 + n$ niekada nėra pilnasis kvadratas jokiam natūraliajam $n$. (Patarimas: naudokitės kvadratų spąstais).
    \item Raskite visus pirminius $p$ ir $q$, kuriems teisinga lygybė $p^2 - q^2 = p+q$.
    \item Raskite visus pirminius skaičius $p$, su kuriais reiškinys $p^2 - 2p$ yra pilnasis kvadratas.
    \item Skaičiai $p$ ir $p^2 + 2$ abu yra pirminiai. Raskite visus tokius galimus $p$. (Patarimas: patikrinkite dalumą iš 3 naudodami $p=3k \pm 1$).
    \item Raskite visus pirminius $p$, su kuriais $3p + 4$ yra pilnasis kvadratas.
    \item Išspręskite lygtį natūraliaisiais skaičiais: $x^2 - y^2 = 2p$, kur $p$ yra nelyginis pirminis skaičius. Jei sprendinių nėra, griežtai įrodykite kodėl. (Patarimas: apsvarstykite $x$ ir $y$ lyginumą).
    \item Raskite visus natūraliuosius $n$, su kuriais skaičius $n^2 + 7$ yra pilnasis kvadratas.
    \item Ar skaičius $p^4 + 4$ gali būti pirminis, jeigu žinoma, kad $p$ pats yra pirminis skaičius?
\end{enumerate}

\subsection*{Sunkūs (Nacionalinių ir aukštesnių olimpiadų lygis)}
\setcounter{enumi}{20}
\begin{enumerate}
    \item Raskite visus pirminius skaičius $p$ ir $q$, tenkinančius diofantinę lygtį $p^2 - 2q^2 = 1$.
    \item Raskite visus pirminius skaičius $p$, kuriems reiškinys $p^3 - p + 1$ yra pilnasis kvadratas.
    \item Raskite visus pirminius $p$, kuriems esant skaičius $8p^2 + 1$ irgi yra pirminis. Įrodykite, kad kitų sprendinių nėra.
    \item Raskite visus natūraliuosius skaičius $n$, kuriems esant $n^4 + n^2 + 1$ yra pirminis skaičius. (Patarimas: pridėkite ir atimkite $n^2$, pritaikykite Sophie Germain idėją).
    \item (Fermat) Įrodykite, kad lygtis $x^3 - y^3 = p$ turi sprendinių natūraliaisiais skaičiais $(x, y)$ tik tokiems pirminiams skaičiams $p$, kurie užrašomi forma $p = 3k^2 + 3k + 1$.
    \item Raskite visus pirminius skaičius $p, q$, su kuriais suma $p^2 + q^2$ yra pilnasis kvadratas. Jei tokių nėra, įrodykite.
    \item Raskite visus pirminius skaičius $p$, kuriems esant lygtis $p = x^4 + 4y^4$ turi sprendinių natūraliaisiais skaičiais $x, y$.
    \item Įrodykite, kad jei $p \ge 5$ yra pirminis skaičius, tai reiškinys $p^2 - 1$ visada tobulai dalijasi iš $24$. (Patarimas: naudokite faktą, kad $p$ nesidalija nei iš $2$, nei iš $3$).
    \item Raskite visus pirminius $p$, su kuriais reiškinys $2^p + p^2$ yra pirminis skaičius.
    \item (Iššūkis) Raskite visus pirminius skaičius $p$, kuriems esant trupmena $\frac{2^{p-1}-1}{p}$ yra sveikasis skaičius ir tobulas pilnasis kvadratas.
\end{enumerate}

\end{document}
